\begin{abstract}
Error-driven world-model adaptation is reactive: the transition that reveals a dynamics mismatch has already occurred. In striking, throwing, and other single-contact tasks, its feedback can arrive after the consequential action is committed. We formulate this first-contact limitation and introduce \method{} (Context-Indexed Residual Adapters), which retrieves previously learned dynamics corrections before interaction. A frozen visual world model is augmented with persistent probabilistic memories, a context posterior combining cues and dynamics evidence, and a factor-structured prior that transfers knowledge to unseen combinations. Memory updates use context probabilities fixed before the current prediction error, permitting a uniform error bound that separates estimation noise, prior mismatch, and contamination by other contexts. The analysis connects memory accuracy and context log loss to first-contact prediction and decision error, and characterizes the observation graph required to identify an unseen additive combination. The same contextual belief supports decision-aware probing and reuse of safety margins: information has zero decision value when all plausible models share an optimal action, while a per-memory calibration rule bounds long-run latent barrier failures under explicit enforcement and boundedness assumptions. A counterfactual decomposition separates improvement through memory selection from improvement through parameter learning. Controlled evaluation protocols isolate recurrence, cue reliability, compositional transfer, calibration, and probing, with the first consequential prediction as the primary endpoint. The resulting framework treats adaptation as inference over a reusable physical repertoire, with new learning reserved for uncertainty that experience has not already resolved.
\end{abstract}
